\section{Results}

\subsection{Aggregate TFP ended below its 2004 level}

Figure~\ref{fig:paper_total_index} reports the chained TFP index for the total economy. The index increased from 100 in 2004 to 100.1 in 2005 and 102.3 in 2006, then declined and remained below 100 from 2009 onward. It reached 98.5 in 2024. The cumulative decline was 1.5 percent, equivalent to an annualized log growth rate of $-0.073$ percent.

\begin{figure}[H]
  \centering
  \includegraphics[width=0.93\textwidth]{paper_fig_ptf_total_index.png}
  \caption{Chained TFP index for the total economy, 2004--2024}
  \label{fig:paper_total_index}
\end{figure}

Most of the aggregate decline occurred before 2011. Aggregate TFP averaged $-0.235$ percent per year in 2005--2010. The subsequent averages were close to zero: 0.022 percent in 2011--2015, $-0.040$ in 2016--2019, and $-0.001$ in 2020--2024. The early deterioration stopped, but aggregate TFP did not enter a sustained growth path.

\subsection{Aggregate stagnation concealed sectoral paths moving in opposite directions}

Figure~\ref{fig:paper_heatmap} displays the 20 annual TFP growth rates that determine the long-run results. The activities are ranked from the lowest to the highest cumulative TFP growth, with the total economy shown in the final row. Mining recorded negative TFP growth in 17 of 20 years and manufacturing in 15. Agriculture and trade recorded positive TFP growth in 15 years; social services did so in 14.

\begin{figure}[H]
  \centering
  \includegraphics[width=\textwidth]{paper_fig_ptf_heatmap.png}
  \caption{Annual TFP growth for the total economy and nine activities, 2005--2024}
  \label{fig:paper_heatmap}
\end{figure}

The annual observations also show why long-run averages should not be interpreted as smooth trends. Construction TFP fell by 9.8 percentage points in 2020. Trade TFP increased by 6.4 points in 2021, while finance and real estate fell by 6.7 points. These changes affect the full-period averages, but the sign frequencies show that the strongest long-run results are not produced by a single year alone.

Figure~\ref{fig:paper_sector_performance} reports the cumulative and annualized measures. Trade had the highest cumulative TFP growth, 23.2 percent, equivalent to an annualized rate of 1.04 percent. Agriculture increased by 21.8 percent and social services by 14.2 percent. Mining had the largest decline, 32.1 percent, equivalent to an annualized rate of $-1.93$ percent. Finance and real estate lost 19.0 percent and construction 11.1 percent.

\begin{figure}[H]
  \centering
  \includegraphics[width=\textwidth]{paper_fig_ptf_sector_performance.png}
  \caption{Cumulative TFP growth from 2004 to 2024 and annualized growth over 2005--2024}
  \label{fig:paper_sector_performance}
\end{figure}

\subsection{Positive and negative sectoral contributions nearly offset}

Sectoral TFP growth does not determine an activity's aggregate contribution on its own. The contribution also depends on its annual aggregation weight. Figure~\ref{fig:paper_contributions} and Table~\ref{tab:paper_long_run} report both dimensions. Four activities contributed 6.308 cumulative log points, or 0.315 percentage points per year. Five activities subtracted 7.770 cumulative log points, or 0.389 points per year. The net contribution was $-1.462$ log points, equivalent to the aggregate annualized TFP growth rate of $-0.073$ percent.

\begin{figure}[H]
  \centering
  \includegraphics[width=\textwidth]{paper_fig_ptf_contributions.png}
  \caption{Activity contributions to aggregate TFP, 2005--2024}
  \label{fig:paper_contributions}
\end{figure}

\begin{table}[H]
\centering
\caption{Sectoral TFP performance and contribution to the total economy, 2005--2024}
\label{tab:paper_long_run}
\footnotesize
\setlength{\tabcolsep}{3pt}
\begin{tabular}{p{4.3cm}rrrr}
\toprule
Activity & Average weight & Cumulative TFP & Annualized TFP & Contribution \\
 & (\%) & growth, 2004--24 (\%) & growth (\%) & (pp per year) \\
\midrule
Trade, hotels, and restaurants & 12.98 & 23.17 & 1.04 & 0.130 \\
Social services & 16.05 & 14.16 & 0.66 & 0.112 \\
Agriculture & 6.11 & 21.84 & 0.99 & 0.056 \\
Transport and communications & 8.63 & 4.26 & 0.21 & 0.017 \\
Electricity, gas, and water & 3.42 & -5.53 & -0.28 & -0.008 \\
Construction & 8.64 & -11.06 & -0.59 & -0.054 \\
Manufacturing & 23.43 & -5.15 & -0.26 & -0.058 \\
Mining & 6.08 & -32.09 & -1.93 & -0.115 \\
Finance and real estate & 14.66 & -19.01 & -1.05 & -0.154 \\
\midrule
\textbf{Total economy} & \textbf{100.00} & \textbf{-1.45} & \textbf{-0.07} & \textbf{-0.073} \\
\bottomrule
\end{tabular}
\vspace{0.15cm}
\begin{minipage}{0.96\textwidth}
\footnotesize \textit{Notes:} Cumulative growth compares the 2024 chained index with 2004. Annualized TFP is the mean of the 20 annual log growth rates from 2005 through 2024. Sectoral cumulative and annualized TFP describe performance within each activity and do not add across activities. The contribution multiplies annual activity TFP by its annual weight before averaging; the nine contributions add to aggregate TFP. Figures may not add because of rounding.
\par\textit{Source:} Authors' calculations based on DANE, 2025 TFP annex.
\end{minipage}
\end{table}


Trade made the largest positive contribution, at 0.130 percentage points per year. Its average weight was 13.0 percent and its own TFP increased at an annualized rate of 1.04 percent. Social services contributed 0.112 points. Agriculture had a lower average weight of 6.1 percent, so its contribution was 0.056 points. Transport contributed 0.017 points.

Finance and real estate made the largest negative contribution, at $-0.154$ percentage points per year. Its own TFP fell at an annualized rate of 1.05 percent and its average weight was 14.7 percent. Mining subtracted 0.115 points. Its average weight was only 6.1 percent, but its own TFP decline was much larger at 1.93 percent per year.

Manufacturing and construction subtracted 0.058 and 0.054 percentage points per year, respectively. Manufacturing had the largest average weight, 23.4 percent, but its own TFP fell at an annualized rate of 0.26 percent. Construction had an average weight of 8.6 percent and a larger sectoral decline of 0.59 percent. These cases show why rankings based only on sectoral TFP growth do not identify the activities with the largest aggregate effect.

\subsection{A zero-TFP counterfactual raises production-account output growth}

Figure~\ref{fig:paper_counterfactual} shows the accounting counterfactual. Production-account output grew by 3.26 percent per year between 2005 and 2024. Setting annual TFP growth to zero in the four weakest activities raises the average to 3.59 percent. The difference is 0.331 percentage points per year.

\begin{figure}[H]
  \centering
  \includegraphics[width=0.95\textwidth]{paper_fig_ptf_counterfactual.png}
  \caption{Observed and counterfactual production-account output growth, 2005--2024}
  \label{fig:paper_counterfactual}
\end{figure}

Finance and real estate and mining account for 0.269 of the 0.331 additional points. Setting finance and real estate TFP growth to zero adds 0.154 points per year, and doing the same for mining adds 0.115. Construction adds 0.054 points and electricity, gas, and water adds 0.008. These changes equal the negative of their average TFP contributions because the exercise holds all other components fixed.

The counterfactual raises output growth in 15 of the 20 years. It lowers growth in 2005, 2006, 2009, 2022, and 2024 because the four selected activities made a positive joint TFP contribution in those years. Compounding the annual differences implies a production-account output level in 2024 that is 6.8 percent above the observed level.

This result should not be interpreted as a policy estimate. It shows the amount of observed growth associated with the negative TFP contributions under the KLEMS identity. It does not show that zero TFP growth was feasible, which intervention could have achieved it, or how factor use and relative prices would have responded.
